\subsection{Validity of Geometric Descriptions}
  \label{subsec:validity-of-geometry}

  The emergence of curvature and metric structure does not imply
  universal validity of geometric descriptions.
  Geometry functions as an effective language,
  applicable only when the relational substrate admits
  a smooth and locally stable continuum approximation.

  Within such regimes, geometric relations are robust:
  local curvature, geodesic deviation, and horizon structure
  depend weakly on microscopic relational details,
  explaining the observed universality of geometric behavior.

  When injectivity or smoothness conditions fail,
  spacetime ceases to be operationally meaningful.
  This signals the breakdown of the geometric description,
  not of the underlying relational structure.

  Geometry therefore acts as an internal consistency condition
  for emergent spacetime descriptions,
  rather than as a fundamental law.

  \paragraph{Relational encoding of mass.}

    The mass parameter in the Schwarzschild-type geometry
    does not arise from a fundamental source term.
    It reflects a localized and persistent modification
    of relational connectivity.

    Operationally, mass corresponds to a reduction in
    admissible relaxation paths,
    modifying the local spectral response of the relational Laplacian
    while preserving global admissibility.

    In the projectable regime,
    this alteration appears as a radial deformation of effective distances,
    whose asymptotic form reproduces the Schwarzschild mass parameter.

    Mass thus quantifies the integrated obstruction to relational relaxation,
    not the presence of an external source.
